\section{Limitations and Future Work}
\label{sec:limitations}

\paragraph{Decode throughput (primary).}
The fused kernel is expressed as many small ops in a compiled graph and does not yet match native
fused attention (\S\ref{sec:analysis}); the exact-decode ablation localizes the gap to the kernel
dispatch, not the algorithm. A hand-written fused decode---a custom Metal op, or a CUDA/Triton
kernel on NVIDIA hardware---is the natural path to closing it. We have deliberately structured the
throughput figures and the main results table so that a future \textbf{CUDA/Triton fused-decode}
measurement can be added as one more series without restructuring the evaluation; no such number is
reported or estimated here, and the placeholder is labelled as future work. Until then, the
adaptive policy keeps the cost off the short-context path entirely.

\paragraph{Compression ratio and the residual memory ceiling.}
Exact recall is bought with memory: at the default budget a quarter of every block is stored
verbatim, so the per-block reduction is $\approx\!2.85\times$, not the order of magnitude a pure
low-rank scheme suggests (\S\ref{sec:compression}). The residuals also set a scaling ceiling---at a
million tokens the residual K/V alone would dominate the pool---so very long contexts need either a
smaller budget (trading recall, per E6) or a two-stage router that gathers residuals only for
selected blocks (which the runtime already does at decode, but the store still holds them). A
budget that adapts to per-block content, rather than a fixed $\maxres$, is the natural refinement.

\paragraph{Router cost.}
The default residual-key router scores every block's anchor and top residual keys each step, an
$O(nR\dhead)$ cost linear in context. It is cheap relative to the kernel at the scales we test, but
at extreme context a cheap pre-filter feeding the exact rerank would be needed---taking care that
the pre-filter cannot drop the needle's block, the failure mode that motivated the exact router in
the first place.

\paragraph{Prefill cost and peak.}
Prefill remains exact and therefore $O(N)$ in attention, with the per-block SVD and residual
extraction added on top, making it the slowest phase and the dominant wall-time cost at large
context. The global memory peak is also set during prefill (\S\ref{sec:analysis}); \diffkv{} does
not reduce it. Compressed or sparse prefill is a known direction (prototyped in the repository) but
is out of scope here.

\paragraph{Evaluation breadth.}
We report one model (Qwen2.5-1.5B int4), one device (Apple~M3, $8.6$\,GB), and one probe
(needle-in-a-haystack) with greedy decoding. Exact-string recall is a sharp but narrow correctness
signal; broader claims require additional models and sizes, perplexity and downstream-quality
metrics, multi-needle and varied-depth probes, and batched/multi-session throughput. The serving
stack supports batching and multi-session residency; we did not exercise them here.

\paragraph{The grounding module.}
The semantic-retrieval/factual-binding layer is inactive in the benchmarked runtime
(\S\ref{sec:impl}). Whether it improves grounded long-context QA on top of the compressed cache
is an open, separately evaluable question.

\paragraph{Portability.}
The runtime is specialized to MLX and to the Qwen2 attention module it patches. A portable
PyTorch/Triton backend and a C++/GGML reimplementation target the same design in the repository but
are under active development and are not evaluated here.
